% =====================================================================
%  CAPÍTULO 2 — ESTADO DA ARTE / REVISÃO BIBLIOGRÁFICA
% =====================================================================
\chapter{Estado da Arte e Enquadramento Teórico}
\label{cap:estado_arte}

Este capítulo divide-se em duas partes. A primeira (\Cref{sec:solucoes})
analisa as soluções existentes no mercado que se aproximam, total ou
parcialmente, dos objetivos deste projeto, identificando o que fazem bem
e o que lhes falta. A segunda (\Cref{sec:enquadramento}) introduz os
conceitos teóricos fundamentais — modelos de linguagem, multimodalidade,
geração de narrativa e geração de vídeo — necessários para compreender as
opções técnicas tomadas nos capítulos seguintes.

\section{Análise de Soluções Existentes}
\label{sec:solucoes}

\subsection{Os ``Big Four'' da genealogia digital}
\label{sub:bigfour}

O mercado da genealogia digital é dominado por quatro grandes plataformas
que se distinguem pela escala dos seus arquivos de registos históricos,
pela extensão das árvores genealógicas colaborativas e pelas ligações
genéticas que oferecem~\cite{familytreemag}:

\begin{itemize}[leftmargin=*]
  \item \textbf{Ancestry.com} — dominante no mercado norte-americano,
  forte integração com testes de \textsc{dna}, particularmente eficaz na
  investigação de ascendência nos EUA.
  \item \textbf{MyHeritage} — mais vocacionada para famílias
  internacionais, reconhecida pelas suas ferramentas de fotografia e
  \textsc{dna}.
  \item \textbf{FamilySearch} — completamente gratuita, mantida pela
  Igreja SUD, com um alcance verdadeiramente global.
  \item \textbf{Findmypast} — especializada em registos do Reino Unido e
  da Irlanda.
\end{itemize}

Estas plataformas são exemplares na agregação e indexação de registos,
mas o seu foco é a \emph{descoberta} de factos genealógicos, não a
\emph{construção de narrativa}.

\subsection{MyHeritage — o caso de estudo mais relevante}
\label{sub:myheritage}

De entre as quatro, a \textbf{MyHeritage} é a plataforma mais próxima
daquilo que este projeto pretende construir, sobretudo pelo seu
investimento em funcionalidades de IA generativa aplicada à memória
familiar~\cite{myheritage_rootstech}:

\begin{itemize}[leftmargin=*]
  \item \textbf{Deep Nostalgia} (2020) — anima fotografias antigas,
  conferindo movimento facial realista a retratos estáticos.
  \item \textbf{DeepStory} (2022) — permite que os antepassados
  ``contem'' as suas próprias histórias de vida, gerando vídeo a partir
  de retratos e de texto.
  \item \textbf{AI Biographer} (2024) — gera, para qualquer antepassado,
  uma entrada de estilo enciclopédico, com contexto histórico,
  fotografias e notas de rodapé que citam a origem dos factos
  genealógicos. Nos dois primeiros meses após o lançamento foram geradas
  cerca de 100.000 biografias~\cite{myheritage_rootstech}.
\end{itemize}

O \emph{AI Biographer} é particularmente relevante por demonstrar a
viabilidade comercial da geração automática de texto biográfico
\emph{ancorado em factos}. Distingue-se, contudo, do presente projeto em
dois aspetos: foca-se em texto (e não em vídeo documental narrado a
partir do arquivo do utilizador) e opera dentro de um ecossistema fechado.

\subsection{Google Photos e Apple Memories — o outro lado}
\label{sub:googlephotos}

Do lado das aplicações de consumo, a IA integrada em serviços como o
\emph{Google~Photos} reanima automaticamente o passado, selecionando
algoritmicamente fotografias que evocam pessoas, locais e momentos. A
funcionalidade \emph{Cinematic Photos and Memories} recorre a
\emph{machine learning} para gerar pequenos clipes animados com um efeito
de profundidade 3D; segundo a Google, o envolvimento dos utilizadores com
a funcionalidade \emph{Memories} cresceu 40\% num ano. Estas soluções são
tecnicamente sofisticadas e visualmente agradáveis, mas o seu resultado é
um \emph{slideshow} estético, não uma \emph{história}.

\subsection{Síntese comparativa e lacuna identificada}
\label{sub:lacuna}

A \Cref{tab:comparativo} sintetiza a análise. O denominador comum às
soluções existentes é a ausência de \textbf{profundidade narrativa} e de
\textbf{integração multifonte}: nenhuma cruza simultaneamente
fotografias, documentos digitalizados, árvores genealógicas e metadados
para construir uma narrativa coerente e factualmente ancorada. É
precisamente essa lacuna que este projeto se propõe preencher.

\begin{table}[H]
\centering
\caption{Análise comparativa das soluções existentes face aos objetivos
do projeto.}
\label{tab:comparativo}
\small
\begin{tabularx}{\linewidth}{@{}p{3.1cm}YYcc@{}}
\toprule
\textbf{Solução} & \textbf{Foco principal} & \textbf{Integra
multifonte?} & \textbf{Narrativa} & \textbf{Vídeo} \\
\midrule
Ancestry / FamilySearch / Findmypast & Registos e árvores genealógicas &
Parcial & Não & Não \\
MyHeritage (AI Biographer) & Biografia textual ancorada & Sim
(genealogia) & Sim (texto) & Parcial \\
Google Photos / Apple Memories & \emph{Slideshows} automáticos & Não &
Não & Sim (estético) \\
\textbf{Este projeto} & Narrativa multimédia do arquivo & \textbf{Sim} &
\textbf{Sim} & \textbf{Sim (documental)} \\
\bottomrule
\end{tabularx}
\end{table}

\section{Enquadramento Teórico}
\label{sec:enquadramento}

\subsection{Modelos de Linguagem de Grande Escala (LLM)}
\label{sub:llm}

Um \emph{Large Language Model} (\textsc{llm}) é, na sua essência, um
modelo treinado para prever a palavra (ou \emph{token}) seguinte numa
sequência de texto. Treinado sobre quantidades massivas de texto —
páginas da internet, livros, artigos científicos — o modelo aprende
padrões estatísticos da linguagem ao ponto de conseguir gerar texto novo,
coerente e contextualmente adequado~\cite{brown2020gpt3}. O mecanismo
central que tornou estes modelos viáveis é a \emph{self-attention},
introduzida na arquitetura \emph{Transformer}~\cite{vaswani2017attention},
que permite ao modelo ponderar simultaneamente as relações entre todas as
palavras de uma sequência, atribuindo pesos diferentes a cada par.

A forma de comunicar com um \textsc{llm} é o \emph{prompt} — a instrução
fornecida ao modelo. A qualidade do resultado depende fortemente da
qualidade do \emph{prompt}, uma constatação que deu origem a uma área de
investigação própria, o \textbf{\emph{prompt engineering}}. A diferença é
ilustrativa: um \emph{prompt} fraco como \emph{``escreve sobre uma
família''} produz texto genérico; um \emph{prompt} bem construído, que
forneça factos concretos, defina a perspetiva, o tom, o formato e a
estrutura, produz resultados radicalmente superiores. Esta observação é
central para o \Cref{cap:implementacao}, onde se descreve o sistema de
\emph{templates} do módulo de geração narrativa.

\subsection{Modelos multimodais — ver além do texto}
\label{sub:multimodal}

Os modelos modernos já não processam apenas texto. Os \emph{Multimodal
Large Language Models} (\textsc{mllm}) recebem como entrada texto, imagem
e até áudio e vídeo em simultâneo. Quando uma fotografia é apresentada a
um modelo como o \emph{GPT-4V} ou o \emph{Gemini}, ocorre, em termos
simplificados, o seguinte~\cite{yin2023mllmsurvey}:

\begin{enumerate}[leftmargin=*]
  \item A imagem é dividida em pequenos blocos (\emph{patches}).
  \item Cada \emph{patch} é convertido numa representação numérica
  (\emph{embedding} visual).
  \item Esses \emph{embeddings} visuais são projetados para o mesmo
  espaço dos \emph{embeddings} de texto.
  \item O modelo processa texto e imagem conjuntamente, como se fossem a
  mesma modalidade.
\end{enumerate}

Representado por modelos como o \emph{GPT-4V}, o paradigma \textsc{mllm}
tornou-se um foco crescente de investigação, usando \textsc{llm}
poderosos como ``cérebro'' para executar tarefas multimodais. Capacidades
emergentes surpreendentes — como \emph{escrever histórias a partir de
imagens} — sugerem um caminho promissor~\cite{yin2023mllmsurvey}. É
exatamente esta capacidade que o módulo de ingestão (M1) explora: um
modelo de visão descreve cada fotografia, gerando matéria-prima factual
para a posterior geração narrativa.

\subsection{Geração automática de narrativas}
\label{sub:geracao_narrativa}

Uma narrativa não é uma mera sequência de factos. Uma boa história possui
\emph{estrutura} (início, desenvolvimento, clímax, conclusão),
\emph{perspetiva} (quem narra e de que ponto de vista), \emph{coerência}
(os factos não se contradizem), \emph{arco emocional} (tensão e resolução)
e \emph{contexto} (os eventos ligam-se a um mundo mais amplo). O desafio
da geração automática consiste em obter tudo isto a partir de dados
brutos e desorganizados.

Identificam-se três paradigmas principais de geração de narrativas:

\begin{description}[leftmargin=*]
  \item[Paradigma 1 — \emph{Template-Based}.] O sistema mais simples:
  definem-se modelos fixos e preenchem-se os espaços com dados. Para este
  projeto, os \emph{templates} são o ponto de partida da arquitetura —
  definem a estrutura narrativa (cronológica, por personagem, temática) e
  o \textsc{llm} preenche-a com linguagem natural rica.
  \item[Paradigma 2 — \emph{Retrieval-Augmented Generation} (RAG).]
  Combina uma fase de \emph{retrieval} (busca de informação relevante numa
  base de dados de factos) com uma fase de \emph{generation} (o
  \textsc{llm} usa essa informação para gerar texto)~\cite{lewis2020rag}.
  \item[Paradigma 3 — \emph{Fine-tuning} e modelos especializados.]
  Parte-se de um modelo base (por exemplo o LLaMA~\cite{touvron2023llama})
  e treina-se especificamente em histórias familiares, para que ``fale''
  com o estilo de um biógrafo. É a abordagem mais poderosa, mas também a
  mais exigente em dados e computação.
\end{description}

Para o âmbito deste projeto, o \textbf{RAG é a abordagem mais
pragmática}: não exige treino de modelos próprios, funciona bem tanto com
modelos locais como com APIs existentes, e — crucialmente — ancora a
geração em factos verificados, mitigando a alucinação. O
\Cref{sub:desafios} aprofunda esta justificação.

\subsection{Desafios técnicos da geração narrativa}
\label{sub:desafios}

\begin{description}[leftmargin=*]
  \item[Alucinação.] Os \textsc{llm} podem inventar factos inexistentes;
  na ausência de informação sobre um período, tendem a ``preencher'' com
  eventos plausíveis mas falsos. \emph{Mitigação:} RAG com factos
  verificados — instruir o modelo a usar \emph{apenas} a informação
  fornecida e a admitir a ausência de registo.
  \item[Coerência temporal.] Dados de décadas podem conter
  inconsistências (uma pessoa que aparece numa fotografia datada antes do
  seu nascimento). \emph{Mitigação:} um módulo de validação que verifica
  a consistência da linha temporal antes da geração (papel do M2).
  \item[Identificação de pessoas.] Sem saber quem é quem nas fotografias,
  a narrativa torna-se genérica. \emph{Mitigação:} reconhecimento facial
  (e.g.\ \emph{DeepFace}) combinado com confirmação manual do utilizador.
  Esta funcionalidade situa-se fora do âmbito da versão atual e é
  discutida como trabalho futuro (\Cref{cap:conclusoes}).
  \item[Avaliação de qualidade.] Determinar se uma narrativa gerada é
  ``boa'' exige métricas de avaliação — objetivas (e.g.\ aderência
  linguística) e subjetivas (painéis humanos) — discutidas no
  \Cref{cap:testes}.
\end{description}

\subsection{Geração de vídeo: pipeline programático vs.\ IA de vídeo}
\label{sub:geracao_video}

Existem duas abordagens fundamentalmente distintas para transformar
fotografias num vídeo:

\begin{description}[leftmargin=*]
  \item[Abordagem A — \emph{pipeline} programático.] Monta-se o vídeo
  como faria um editor: define-se a duração de cada fotografia,
  adiciona-se narração por \textsc{tts}, música de fundo e o efeito
  \emph{Ken~Burns} (zoom/panorâmica lentos que dão vida a imagens
  estáticas), exportando um ficheiro \textsc{mp4}. É o caminho com
  bibliotecas de código aberto como o \emph{MoviePy}~\cite{moviepy} e o
  \emph{FFmpeg}~\cite{ffmpeg}, de custo nulo e plenamente controlável.
  \item[Abordagem B — geração de vídeo por IA.] Recorre a APIs como o
  \emph{Runway}~\cite{runway}, \emph{Pika}, \emph{Luma} ou \emph{Kling} para, a partir
  de uma fotografia estática, gerar segundos de vídeo animado com
  movimento realista. É mais impressionante, mas também mais cara e
  complexa, e levanta questões de consistência e de custo por geração.
\end{description}

Para um projeto académico, um \emph{slideshow} narrado de qualidade
construído pela Abordagem A é \textbf{suficiente e academicamente
defensável}, cumprindo os objetivos sem dependência de APIs dispendiosas.
Foi essa a opção adotada, conforme detalhado no \Cref{cap:implementacao}.
A \Cref{tab:videoapis} resume, para referência, as principais APIs de
geração de vídeo da Abordagem B.

\begin{table}[H]
\centering
\caption{Principais APIs de geração de vídeo por IA (Abordagem B), para
referência. Não utilizadas na versão atual por motivos de custo.}
\label{tab:videoapis}
\small
\begin{tabularx}{\linewidth}{@{}lYl@{}}
\toprule
\textbf{Ferramenta} & \textbf{Ideal para} & \textbf{API} \\
\midrule
Pika Labs            & Prototipagem rápida                  & Sim \\
Runway Gen-4         & Qualidade profissional               & Sim \\
Luma (Dream Machine) & Conversão foto $\rightarrow$ vídeo realista & Sim \\
Kling AI             & Alternativa de boa qualidade         & Sim \\
\bottomrule
\end{tabularx}
\end{table}

\subsection{Síntese de voz (Text-to-Speech)}
\label{sub:tts}

A síntese de voz (\textsc{tts}) converte o texto narrativo em fala,
dando ``voz'' à história. A qualidade da voz afeta diretamente o impacto
emocional do vídeo: uma narração natural e expressiva tem um efeito
completamente distinto de uma voz robótica. Entre as opções de código
aberto/gratuitas consideradas: o \emph{gTTS} (simples, dependente de
internet, qualidade razoável mas pouco expressiva), as vozes neuronais do
\emph{edge-tts} (qualidade elevada, vozes europeias) e o \emph{Coqui~TTS}
(local, com a possibilidade de \emph{voice cloning} a partir de uma
amostra real — funcionalidade tecnicamente e eticamente rica, mas
exigente em recursos). As opções efetivamente integradas no projeto
discutem-se no \Cref{cap:implementacao}.

\subsection{Ferramentas de código aberto por módulo}
\label{sub:ferramentas}

Um princípio orientador do projeto é o custo essencialmente nulo, através
de ferramentas de código aberto e camadas gratuitas. A
\Cref{tab:ferramentas} resume, por módulo, as tecnologias consideradas no
estudo inicial — várias delas efetivamente adotadas.

\begin{table}[H]
\centering
\caption{Ferramentas de código aberto / gratuitas consideradas, por
módulo.}
\label{tab:ferramentas}
\small
\begin{tabularx}{\linewidth}{@{}p{3.4cm}Y@{}}
\toprule
\textbf{Módulo} & \textbf{Ferramentas} \\
\midrule
M1 — Ingestão & Google Gemini (visão, camada gratuita); Llama 3.2 Vision
via Ollama (local); extração \textsc{exif}; Tesseract \textsc{ocr};
\emph{parser} GEDCOM. \\
M2 — Temporal & Lógica própria em Python; NetworkX para o grafo de
parentesco. \\
M3 — Narrativa & Groq (Llama 3.3 70B, camada gratuita na nuvem); Ollama +
Llama 3.2 / Mistral / Gemma (local); Google Gemini Flash (camada gratuita);
Hugging Face Inference (com limites). \\
M4 — Multimédia & MoviePy; gTTS / edge-tts / Coqui TTS; FFmpeg. \\
\bottomrule
\end{tabularx}
\end{table}

\noindent Onde poderiam surgir custos — APIs comerciais de \textsc{llm}
proprietárias em produção, ou APIs de geração de vídeo (Runway, Pika) —
o projeto evita-os deliberadamente: modelos de camada gratuita (\emph{Groq}
na nuvem e \emph{Ollama} em local) para o texto e um \emph{slideshow}
programático com MoviePy no lugar da geração de vídeo por IA.
